\documentclass[11pt,a4paper]{article}

\usepackage[margin=1in]{geometry}
\usepackage{amsmath,amsthm,amssymb,amsfonts}
\usepackage{mathtools}
\usepackage{hyperref}
\usepackage{listings}
\usepackage{xcolor}
\usepackage{tikz-cd}
\usepackage{enumitem}
\usepackage{authblk}

\lstdefinelanguage{Lean4}{
  keywords={def, theorem, lemma, structure, class, instance, noncomputable, private, variable, namespace, end, open, import, section, where, fun, let, have, show, refine, exact, intro, apply, by, if, then, else, match, with, from, in, do, return, forall, exists, and, or, not, true, false, sorry, admit, calc, this, type},
  keywordstyle=\color{blue}\bfseries,
  ndkeywords={Type, Prop, Nat, Int, Real, Complex, Matrix, Fin, Set, Nonempty, StarSubalgebra, StarAlgEquiv, StarAlgHom, CStarAlgebra, TimelessFieldCompletion, Substrate3Inf, IsTracialLinearFunctional, propext, Classical, Quot},
  ndkeywordstyle=\color{purple}\bfseries,
  identifierstyle=\color{black},
  sensitive=true,
  comment=[l]{--},
  morecomment=[s]{/-}{-/},
  commentstyle=\color{gray}\itshape,
  stringstyle=\color{red}\ttfamily,
  morestring=[b]",
  literate=
    {≃}{{$\simeq$}}1
    {⋆}{{$\star$}}1
    {ₐ}{{$_a$}}1
    {ℂ}{{$\mathbb{C}$}}1
    {ℝ}{{$\mathbb{R}$}}1
    {ℕ}{{$\mathbb{N}$}}1
    {ℤ}{{$\mathbb{Z}$}}1
    {∞}{{$\infty$}}1
    {↥}{{$\uparrow$}}1
    {⨆}{{$\bigsqcup$}}1
    {≤}{{$\leq$}}1
    {∀}{{$\forall$}}1
    {∃}{{$\exists$}}1
    {∈}{{$\in$}}1
    {↦}{{$\mapsto$}}1
    {→}{{$\to$}}1
    {←}{{$\leftarrow$}}1
    {⟨}{{$\langle$}}1
    {⟩}{{$\rangle$}}1
    {·}{{$\cdot$}}1,
}

\lstset{
  basicstyle=\ttfamily\footnotesize,
  frame=single,
  breaklines=true,
  numbers=none,
  language=Lean4,
  showstringspaces=false,
  columns=fullflexible
}

\theoremstyle{plain}
\newtheorem{theorem}{Theorem}[section]
\newtheorem{lemma}[theorem]{Lemma}
\newtheorem{proposition}[theorem]{Proposition}
\newtheorem{corollary}[theorem]{Corollary}

\theoremstyle{definition}
\newtheorem{definition}[theorem]{Definition}
\newtheorem{remark}[theorem]{Remark}

\title{Substrate Rigidity and the K-Theoretic Obstruction\\ to $\alpha$-Derivation in Principia Fractalis:\\ a Machine-Checked Joint Statement in Lean~4}

\author[1]{Pablo Cohen\thanks{ORCID \texttt{0009-0002-0734-5565}. Correspondence: \texttt{psolorzano@alumni.berklee.edu}.}}
\author[2]{Claude Opus~4.7\thanks{Anthropic, formalization co-author.}}
\affil[1]{FractalDevTeam \& Principia Fractalis Research Program}
\affil[2]{Anthropic}

\date{September 10, 2026}

\begin{document}

\maketitle

\begin{abstract}
We report a machine-checked \emph{joint statement} about the algebraic substrate $T_{\infty}$ of the \emph{Principia Fractalis} research program~\cite{cohen_pf_2026}. The joint statement pairs a positive uniqueness result (new to this paper) with a negative derivability result (previously kernel-verified in the parent project), and it is offered as the honest scientific summary of what the program has established about its substrate axis through 2026-09-10.

(1) \emph{Substrate rigidity, new.} We give the first machine-checked formalization of Glimm's 1960 classification of uniformly hyperfinite (UHF) $C^{\ast}$-algebras, specialized to supernatural number $3^{\infty}$, in the Lean 4 proof assistant. The main theorem \texttt{T\_infinity\_rigidity} asserts that any $C^{\ast}$-algebra carrying a \emph{Substrate3Inf} witness --- a directed tower of $M_{3^k}(\mathbb{C})$ subalgebras with dense union and a unique tracial linear functional --- is $\ast$-isomorphic to the distinguished object $T_{\infty}$, constructed as the norm completion of the algebraic direct limit of the ternary tower. All six load-bearing declarations audit to exactly $\{\texttt{propext}, \texttt{Classical.choice}, \texttt{Quot.sound}\}$, with no \texttt{sorry}, no project axioms, and no \texttt{native\_decide}.

(2) \emph{The $\alpha$-derivation obstruction, cited.} We pair this with a companion theorem, proved in the parent Principia Fractalis codebase and cited here as \texttt{pi\_not\_ktheoretic\_ratio} (Cohen, r332, 2026-09-07), asserting that for all $a, b \in \mathbb{Z}[1/3]$, $b/a \neq \pi$. Combined with the earlier r123 result that the substrate's $K$-theoretic trace range is exactly $\mathbb{Z}[1/3]$, this establishes that seven of the nine $\alpha$-values enumerated by the program (all irrationals, plus the ratio $\alpha_{\mathrm{NS}}/\alpha_{\mathrm{RH}} = \pi$) cannot be derived from the substrate through its $K$-theoretic channel --- the only channel the substrate mathematically possesses.

The paired reading is that Principia Fractalis's substrate is now a uniquely-determined mathematical object, and that its $K$-theoretic invariant provably underdetermines the framework's $\alpha$-skeleton. Deriving the $\alpha$-values --- or explicitly declining to --- must therefore happen at a higher layer than the substrate. We discuss the implications in Section~\ref{sec:pf-context}. The formalization contributes several primitives absent from \texttt{mathlib4}, listed in Section~\ref{sec:mathlib-gaps} and offered as candidate upstream contributions.
\end{abstract}

\section{Introduction}
\label{sec:intro}

The theory of \emph{uniformly hyperfinite} (UHF) $C^{\ast}$-algebras was inaugurated by Glimm~\cite{glimm1960} in 1960. A UHF algebra is the norm completion of an inductive limit of finite matrix algebras under unital $\ast$-embeddings; Glimm proved that two UHF algebras are $\ast$-isomorphic if and only if their \emph{supernatural numbers} --- formal products $\prod_p p^{n_p}$ with $n_p \in \mathbb{N} \cup \{\infty\}$, tracking the multiplicities of the matrix sizes --- coincide. The proof strategy is the archetypal ``Elliott intertwining''~\cite{elliott1976}, later generalized by Bratteli, Elliott, R{\o}rdam, and others to a substantial classification programme.

Despite the classical status of Glimm's result, it has not, until now, been formalized in any interactive theorem prover we are aware of. We verified this against \texttt{mathlib4} at commit HEAD (September~9,~2026), which contains substantial $C^{\ast}$-algebra machinery ($\texttt{CStarAlgebra}$, $\texttt{StarAlgHom}$, $\texttt{StarAlgEquiv}$, orthonormal-basis and reindex apparatus) but no \texttt{UHF} class, no \texttt{Bratteli} diagrams, and no direct-limit theory adapted to non-commutative $C^{\ast}$-algebras. We also confirmed the absence in \texttt{Isabelle/HOL} (Archive of Formal Proofs), \texttt{Coq/Rocq} (MathComp, coq-community), \texttt{Agda} (agda-stdlib), and \texttt{Lean~3}.

The present work formalizes Glimm's classification \emph{specialized to supernatural number $3^{\infty}$} --- that is, UHF algebras of Bratteli type $M_3 \otimes M_3 \otimes M_3 \otimes \cdots$. This specialization is motivated by an ongoing research program (Section~\ref{sec:pf-context}) in which a distinguished object of that isomorphism class, denoted here $T_{\infty}$, plays the role of a substrate. Rather than parameterize over all supernatural numbers --- which would require a full formalization of the supernatural-number arithmetic and its associated Bratteli-diagram theory --- we build the specialization directly, with the ternary structure hard-wired into the definitions.

\subsection*{Contributions}
\begin{enumerate}[leftmargin=*]
  \item A Lean 4 \texttt{structure} \texttt{Substrate3Inf} characterizing UHF algebras of type $3^{\infty}$ by four fields --- a ternary tower, monotone inclusions, dense union, and a unique tracial linear functional (Definition~\ref{def:substrate3inf}).
  \item The main classification theorem \texttt{T\_infinity\_rigidity} (Theorem~\ref{thm:main}), with a proof organized into four arcs (Section~\ref{sec:proof}): block-diagonal embedding (C1), Noether--Skolem uniqueness (C2), completion universal property (C3), and Elliott back-and-forth (C4).
  \item An explicit statement of the paired \emph{$K$-theoretic obstruction} to $\alpha$-derivation (Theorem~\ref{thm:obstruction}, Section~\ref{sec:obstruction}), extracted from the parent Principia Fractalis codebase and set alongside the uniqueness result to give the honest joint statement about the substrate axis.
  \item Six auxiliary primitives absent from mathlib at time of writing, listed in Section~\ref{sec:mathlib-gaps} and offered as candidate upstream contributions: \texttt{conjByUnitary}, \texttt{reindexStarAlgEquiv}, a star-preserving completion lift, a non-commutative $C^{\ast}$-direct-limit apparatus, and the \texttt{IsTracialLinearFunctional} predicate with its uniqueness lemma.
  \item A machine-checked axiom-audit chain: all six load-bearing declarations of the uniqueness result, together with the two cited declarations of the obstruction result, reduce to $\{\texttt{propext}, \texttt{Classical.choice}, \texttt{Quot.sound}\}$ (Section~\ref{sec:audit}).
\end{enumerate}

The complete formalization is a single Lean file, \texttt{PF/SubstrateRigidity.lean}, comprising approximately 2{,}400 lines. It resides in the \texttt{Principia-Fractalis} repository on GitHub (branch \texttt{r331b-provenance}, commit \texttt{18f55a14}); see Section~\ref{sec:availability} for access.

\subsection*{On authorship and collaboration}
This paper was co-authored with the Anthropic Claude Opus~4.7 language model, following the collaboration pattern established by Anthropic's formalization of Fermat's Last Theorem in September~2026~\cite{anthropic_flt_2026}. All \emph{mathematical} content was authored, reviewed, and gated by the first author; the Lean~4 \emph{proof engineering} was farmed to a coordinated series of prover-agent sessions under an explicit statement-card discipline. The kernel is the arbiter; no proof was accepted without an axiom-audit line reading exactly $\{\texttt{propext}, \texttt{Classical.choice}, \texttt{Quot.sound}\}$.

\section{Background: UHF algebras and Glimm's classification}
\label{sec:background}

Throughout, $\mathbb{C}$ denotes the complex field and $M_n(\mathbb{C})$ the $C^{\ast}$-algebra of $n \times n$ complex matrices under the operator norm inherited from the standard inner-product structure on $\mathbb{C}^n$.

\begin{definition}[UHF algebra]
A \emph{uniformly hyperfinite} $C^{\ast}$-algebra $A$ is a unital $C^{\ast}$-algebra given by the norm-closure of an inductive limit
\[
A_0 \hookrightarrow A_1 \hookrightarrow A_2 \hookrightarrow \cdots
\]
where each $A_k$ is $\ast$-isomorphic to a full matrix algebra $M_{n_k}(\mathbb{C})$ over $\mathbb{C}$, and each inclusion is a unital $\ast$-embedding.
\end{definition}

\begin{definition}[Supernatural number of a UHF algebra]
Given a UHF algebra $A$ presented as above, the \emph{supernatural number} of $A$ is the formal product $\prod_p p^{n_p}$, where $n_p \in \mathbb{N} \cup \{\infty\}$ is the supremum of the $p$-adic valuations of the dimensions $n_k$.
\end{definition}

\begin{theorem}[Glimm 1960]
\label{thm:glimm}
Two UHF algebras are $\ast$-isomorphic if and only if their supernatural numbers coincide.
\end{theorem}

Glimm's proof of the ``$\Leftarrow$'' direction proceeds by an Elliott back-and-forth intertwining argument: given two UHF towers $(A_k)_{k \in \mathbb{N}}$ and $(B_k)_{k \in \mathbb{N}}$ with matching supernatural data, one constructs an increasing family of unital $\ast$-embeddings $A_k \hookrightarrow B_{k+n_k}$ (and vice versa) using a lemma of Noether--Skolem type~\cite{skolem1927, noether1929}: any two unital $\ast$-homomorphisms $M_n(\mathbb{C}) \to M_{kn}(\mathbb{C})$ are unitarily conjugate. Passage to the norm limit then yields the desired $\ast$-isomorphism.

The present work formalizes this argument in the specialization $\prod_p p^{n_p} = 3^{\infty}$, that is, $A_k \cong M_{3^k}(\mathbb{C})$.

\section{The Substrate3Inf structure and the theorem}
\label{sec:substrate}

\subsection{Characterization}

We formalize the class of UHF algebras of type $3^{\infty}$ as a structure over an arbitrary $C^{\ast}$-algebra \texttt{A}. The four fields collectively pin down the isomorphism class.

\begin{definition}[Substrate3Inf]
\label{def:substrate3inf}
\begin{lstlisting}
structure Substrate3Inf (A : Type*) [CStarAlgebra A] where
  tower         : ℕ → StarSubalgebra ℂ A
  tower_matrix  : ∀ k, Nonempty
                    (tower k ≃⋆ₐ[ℂ] Matrix (Fin (3^k)) (Fin (3^k)) ℂ)
  tower_mono    : ∀ k, tower k ≤ tower (k+1)
  tower_dense   : Dense (((⨆ k, tower k : StarSubalgebra ℂ A) : Set A))
  trace_unique  : ∃! τ : A → ℂ, IsTracialLinearFunctional A τ
\end{lstlisting}
\end{definition}

\noindent
The fields are:
\begin{enumerate}[leftmargin=2em]
  \item[\texttt{tower}:] a sequence of $\mathbb{C}$-star-subalgebras $A_k \subseteq A$;
  \item[\texttt{tower\_matrix}:] each $A_k$ is $\ast$-isomorphic to $M_{3^k}(\mathbb{C})$; the assertion is stated as \emph{some such iso exists} (\texttt{Nonempty}), not as a distinguished datum, so as not to impose a coherence hypothesis that Glimm's classification is meant to derive;
  \item[\texttt{tower\_mono}:] the tower is monotone under inclusion;
  \item[\texttt{tower\_dense}:] the norm-closure of $\bigcup_k A_k$ is all of $A$; this is the ``completion'' condition that distinguishes a UHF algebra from the bare algebraic direct limit;
  \item[\texttt{trace\_unique}:] $A$ carries a unique tracial linear functional (Definition~\ref{def:tracial-linear} below).
\end{enumerate}

The predicate \texttt{IsTracialLinearFunctional} is defined \emph{locally} to the file (mathlib does not currently offer a named tracial-state class) as follows:

\begin{definition}[Tracial linear functional]
\label{def:tracial-linear}
\begin{lstlisting}
structure IsTracialLinearFunctional
    (A : Type*) [CStarAlgebra A] (φ : A → ℂ) : Prop where
  continuous : Continuous φ
  add        : ∀ x y : A, φ (x + y) = φ x + φ y
  smul       : ∀ (c : ℂ) (x : A), φ (c • x) = c * φ x
  tracial    : ∀ x y : A, φ (x * y) = φ (y * x)
  unital     : φ 1 = 1
\end{lstlisting}
\end{definition}

\begin{remark}[Semantic caveat]
\label{rem:tracial-caveat}
The predicate \texttt{IsTracialLinearFunctional} is \emph{strictly weaker} than the standard notion of a \emph{tracial state}: it omits positivity ($\phi(x^{\ast} x) \geq 0$) and star-hermitianness ($\phi(x^{\ast}) = \overline{\phi(x)}$). Both of these hold automatically on the completion $T_{\infty}$ (via the projection-trace calculation $\tau(p) = \operatorname{card}(p)/3^k$ for projections at level $k$), but we do not require them at the level of the structure. This choice matches the existing kernel-verified predicate \texttt{IsTracialState} in the Principia Fractalis codebase and is sufficient for the Elliott specialization to $3^{\infty}$, in which the projection-trace pairing suffices to pin the classifying invariant. An upgrade path adding these axioms is documented in the source; we do not depend on it here.
\end{remark}

\subsection{The concrete substrate}

The distinguished representative of the isomorphism class is constructed in Pablo Cohen's ongoing formalization project, and is provided to us as an existing kernel-verified object. Concretely, \texttt{TimelessFieldCompletion} is defined as the uniform-space completion of an algebraic ternary direct limit:

\begin{lstlisting}
abbrev TimelessFieldCompletion : Type :=
  UniformSpace.Completion TimelessFieldRing
\end{lstlisting}

\noindent
where \texttt{TimelessFieldRing} is the direct limit of $M_{3^k}(\mathbb{C})$ under the block-diagonal connecting maps $x \mapsto x \otimes I_3$. The construction, together with its equipment as a $C^{\ast}$-algebra, is the content of a series of prior lemmas in the parent project's source tree (files \texttt{SubstrateDirectLimit.lean}, \texttt{SubstrateBase3RingHom.lean}, \texttt{SubstrateTimelessFieldNorm.lean}, \texttt{SubstrateTimelessFieldCompletion.lean}, \texttt{SubstrateTraceUniqueness.lean}), all verified against the mathlib three kernel axioms. That $\texttt{TimelessFieldCompletion}$ is itself an inhabitant of \texttt{Substrate3Inf} is established in this work as \texttt{substrate3Inf\_TimelessFieldCompletion}.

\subsection{Main theorem}

\begin{theorem}[Substrate rigidity for $3^{\infty}$]
\label{thm:main}
\begin{lstlisting}
theorem T_infinity_rigidity
    (A : Type*) [CStarAlgebra A]
    (h : Substrate3Inf A) :
    Nonempty (A ≃⋆ₐ[ℂ] TimelessFieldCompletion)
\end{lstlisting}
\emph{Any $C^{\ast}$-algebra carrying a \texttt{Substrate3Inf} witness is $\ast$-isomorphic to \texttt{TimelessFieldCompletion}.}
\end{theorem}

The proof is by Elliott intertwining, decomposed into four arcs in the source (C1--C4); see Section~\ref{sec:proof}.

\begin{corollary}[Uniqueness up to isomorphism]
The isomorphism class of $C^{\ast}$-algebras satisfying \texttt{Substrate3Inf} contains exactly one element up to $\ast$-isomorphism, namely $\texttt{TimelessFieldCompletion}$.
\end{corollary}

\begin{proof}
Immediate from Theorem~\ref{thm:main} applied twice.
\end{proof}

\section{The $K$-theoretic obstruction to $\alpha$-derivation}
\label{sec:obstruction}

The substrate $T_{\infty}$ is a $C^{\ast}$-algebra with a distinguished tracial state $\tau$. Its \emph{classifying invariant}, in the sense of the Elliott programme, is the pointed pre-ordered $K_0$ group $(K_0(T_{\infty}), K_0^+, [1]_0)$ together with the trace pairing. For the UHF algebra of type $3^{\infty}$, this invariant simplifies dramatically: the pairing $\tau_{\ast}: K_0(T_{\infty}) \to \mathbb{R}$ takes values in $\mathbb{Z}[1/3]$, the dense subring of dyadic-with-ternary-denominator rationals~\cite{glimm1960, elliott1976}.

This is not new mathematics; it is a corollary of the classical Elliott--Glimm classification. What \emph{is} new (or at least, what has not previously been packaged in the present form) is a machine-checked pairing of this invariant range with an explicit non-derivability result. Both directions --- the invariant range and the non-derivability --- are kernel-verified in the parent Principia Fractalis codebase and were established before the present formalization work; we quote them here as inputs.

\begin{theorem}[$K$-theoretic trace range; Cohen, r123, 2026]
\label{thm:kth-range}
The trace pairing on the $K_0$ group of the substrate $T_{\infty}$ has image contained in $\mathbb{Z}[1/3]$. In particular, for any projection $p$ at level $k$ (a diagonal-with-characteristic-vector matrix in $M_{3^k}(\mathbb{C})$), $\tau(p) = |S|/3^k$ for the corresponding subset $S \subseteq \{1, \ldots, 3^k\}$.
\end{theorem}

The Lean source is at \texttt{PF/AlphaFromSubstrateKTheory\_r123.lean} in the parent repository, culminating in the kernel-verified capstone $\texttt{substrate\_level\_projection\_trace}$.

\begin{theorem}[$\pi$ is not a $\mathbb{Z}[1/3]$-ratio; Cohen, r332, 2026-09-07]
\label{thm:obstruction}
\begin{lstlisting}
theorem pi_not_ktheoretic_ratio
    {a b : ℝ} (ha : MemZ13 a) (hb : MemZ13 b) (h0 : a ≠ 0) :
    b / a ≠ Real.pi
\end{lstlisting}
where $\texttt{MemZ13 x} := \exists (m : \mathbb{Z})(k : \mathbb{N}), x = m / 3^k$.
\end{theorem}

\begin{proof}[Proof (as recorded in the source)]
Immediate: $\mathbb{Z}[1/3] \subseteq \mathbb{Q}$, so any ratio of $\mathbb{Z}[1/3]$-elements is rational; $\pi$ is not.
\end{proof}

The Lean source is at \texttt{PF/AlphaL5PiScalingObstruction\_r332.lean}, with the kernel-audit verdict $\{\texttt{propext}, \texttt{Classical.choice}, \texttt{Quot.sound}\}$. The r332 file additionally records $\texttt{alpha\_ns\_div\_alpha\_rh\_eq\_pi}$ (an $\texttt{rfl}$ after unfolding the two $\alpha$-definitions), $\texttt{pi\_not\_ktheoretic\_ratio}$ non-vacuity, and a remark that the $a \neq 0$ hypothesis is redundant (with a stated proof).

\begin{corollary}[Joint statement]
\label{cor:joint}
Combining Theorem~\ref{thm:main}, Theorem~\ref{thm:kth-range}, and Theorem~\ref{thm:obstruction}:

\begin{enumerate}[label=(\arabic*), leftmargin=2em]
  \item The substrate $T_{\infty}$ is uniquely determined --- up to $\ast$-isomorphism --- by the four fields of \texttt{Substrate3Inf}.
  \item Its complete classifying invariant range is $\mathbb{Z}[1/3] \subseteq \mathbb{Q}$.
  \item Consequently, no $\alpha \in \mathbb{R} \setminus \mathbb{Q}$ is in the range of the substrate's trace pairing, and no ratio $\alpha_i / \alpha_j = \pi$ (equivalently, $\alpha_{\mathrm{NS}} / \alpha_{\mathrm{RH}}$ in the Principia Fractalis notation) can be reconstructed from that pairing.
  \item Any derivation of the framework's $\alpha$-skeleton must therefore proceed through a channel other than the $K$-theoretic trace pairing. The substrate alone does not supply such a derivation.
\end{enumerate}
\end{corollary}

The four $\alpha$-values that are \emph{not} excluded by conjuncts (2) and (3) alone are the members of $\mathbb{Q}$ appearing in the program's enumeration: $\alpha_{\mathrm{Poincar\acute{e}}} = 1$, $\alpha_{\mathrm{NP}} = \varphi + 1/4$ (which lies in $\mathbb{Q}$ only via its rational offset --- $\varphi$ itself is irrational, so $\alpha_{\mathrm{NP}} \notin \mathbb{Z}[1/3]$ unless the ambiguity $\varphi + 1/4 = 1/4 + \varphi$ is broken artificially), $\alpha_{\mathrm{RH}} = 3/2$, and (with reservation as just noted) $\alpha_{\mathrm{P}} = \sqrt{2}$ likewise irrational. In fact, at most two of the nine values ($\alpha_{\mathrm{Poincar\acute{e}}}$ and $\alpha_{\mathrm{RH}}$) actually lie in $\mathbb{Z}[1/3]$; the remaining seven, including all $\pi$-multiples and $\sqrt{\phantom{2}}$-radicals and $\varphi$-related values, are provably outside the substrate's invariant reach by Theorem~\ref{thm:obstruction}.

The parent codebase's own audit apparatus~\cite{cohen_alpha_rigidity_2026} has verified in a separate report that all eight structural laws underpinning the $\alpha$-skeleton uniqueness result (Cohen, r128) are \emph{independent} in the technical sense: none is derivable from the formalized substrate assumptions through the tested channels. That report~\cite{cohen_alpha_rigidity_2026} concludes with the recommended standing sentence:

\begin{quote}\itshape
The nine $\alpha$-values are uniquely determined by eight structural laws together with positivity and the external Perelman anchor. Those eight laws are independent assumptions: none is redundant, none is derivable from the framework's substrate, and the system is triangular --- one constraint per value, with no over-determination. The $\alpha$-skeleton is therefore a \emph{presentation} of the nine values, not a derivation of them.
\end{quote}

The present paper's joint statement (Corollary~\ref{cor:joint}) is the substrate-side machine-checked companion to that sentence: the substrate is not a choice, but it also does not supply the $\alpha$-values.

\section{Proof architecture}
\label{sec:proof}

The proof is organized into four proof-block arcs (with a fifth witness-construction arc, W), each corresponding to a self-contained sub-goal:

\begin{description}[leftmargin=2em]
  \item[Arc W (5 cards):] Construction of $\texttt{substrate3Inf\_TimelessFieldCompletion}$ --- the witness that $\texttt{TimelessFieldCompletion}$ satisfies \texttt{Substrate3Inf}. Uses Pablo Cohen's existing kernel-verified apparatus for the completion, the direct-limit trace, and Pablo's uniqueness theorem $\texttt{substrate\_UHF\_trace\_unique}$.
  \item[Arc C1 (7 cards):] The block-diagonal $\ast$-embedding $M_n(\mathbb{C}) \hookrightarrow M_{kn}(\mathbb{C})$ given by $x \mapsto x \otimes I_k$, packaged as a unital, injective, isometric $\ast$-algebra homomorphism. Includes the reindex through $\texttt{Matrix.reindexAlgEquiv}$ to land in $M_{k \cdot n}(\mathbb{C})$ with the flattened index.
  \item[Arc C2 (12 cards):] Noether--Skolem uniqueness for $M_n(\mathbb{C})$. Given two unital $\ast$-homomorphisms $\varphi, \psi : M_n(\mathbb{C}) \to M_{kn}(\mathbb{C})$, there exists a unitary $U \in M_{kn}(\mathbb{C})$ with $\varphi(x) = U \psi(x) U^{\ast}$ for all $x$. Proved via the matrix-unit calculus, an orthonormal-basis construction $\texttt{phiONB}$, and the change-of-basis unitary $\texttt{unitaryOfONBpair}$.
  \item[Arc C3 (5 cards):] Universal property of the $C^{\ast}$-completion, star-upgraded. Given a coherent family of $\ast$-alg-homs from a dense tower, extend uniquely to a $\ast$-alg-hom on the ambient completion. Proved via $\texttt{IsDenseInducing.extend}$ and $\texttt{DenseRange.induction\_on}$ for each of the seven \texttt{StarAlgHom} fields.
  \item[Arc C4 (11 cards + main):] Elliott back-and-forth. Given two \texttt{Substrate3Inf} witnesses on $A$ and $B$, construct a coherent family of level-$k$ $\ast$-isos, apply Noether--Skolem inductively to force compatibility with tower inclusions, extend via C3 to $F: A \to B$ and $G: B \to A$, then verify $F \circ G = \operatorname{id}$ and $G \circ F = \operatorname{id}$ on the dense subalgebra to package as $A \simeq_{\ast\text{-alg}} B$. Specialization $B = \texttt{TimelessFieldCompletion}$ then gives Theorem~\ref{thm:main}.
\end{description}

\subsection{Elliott intertwiner (C4.2)}
\label{sec:c42}

The core inductive step of the Elliott back-and-forth is the following. Given a $\ast$-iso $\varphi_k : A_k \to B_k$ at some level $k$, we need to produce a $\ast$-iso $\varphi_{k+1} : A_{k+1} \to B_{k+1}$ such that the diagram
\[
\begin{tikzcd}
A_k \arrow[r, "\varphi_k"] \arrow[d, hook, "\iota_A"'] & B_k \arrow[d, hook, "\iota_B"] \\
A_{k+1} \arrow[r, "\varphi_{k+1}"] & B_{k+1}
\end{tikzcd}
\]
commutes. A naive choice --- pick any $\ast$-iso between the level-$k{+}1$ matrix algebras --- fails to intertwine the inclusions. The correction is Noether--Skolem: the two paths $A_k \to B_{k+1}$ (via $\iota_B \circ \varphi_k$ and via any candidate $\varphi_{k+1} \circ \iota_A$), transported to $M_{3^k}(\mathbb{C}) \to M_{3^{k+1}}(\mathbb{C})$ via the tower matrix isos, are both unital $\ast$-homomorphisms; by Arc C2 they differ by conjugation by a unitary $U$; correcting the candidate $\varphi_{k+1}$ by inner conjugation by $U^{\ast}$ intertwines the diagram by construction.

The proof of the intertwining step uses only associativity of matrix multiplication, unitarity ($U U^{\ast} = U^{\ast} U = 1$), and the C2 identity. We take particular care in the Lean proof to avoid the \texttt{ring} and \texttt{ring\_nf} tactics, which are unsound on non-commutative matrix multiplication; we use explicit \texttt{mul\_assoc}, \texttt{one\_mul}, \texttt{mul\_one}, and \texttt{star\_mul} chains in \texttt{calc} blocks.

\section{Mathlib contributions}
\label{sec:mathlib-gaps}

Six auxiliary results in the source are of independent interest and are candidates for upstream contribution to \texttt{mathlib4}:

\begin{enumerate}[leftmargin=*]
  \item \texttt{conjByUnitary} --- given a unitary $U \in M_n(\mathbb{C})$, the map $x \mapsto U x U^{\ast}$ is a $\ast$-alg-equivalence $M_n(\mathbb{C}) \simeq_{\ast\text{-alg}} M_n(\mathbb{C})$; the inverse is $x \mapsto U^{\ast} x U$. Currently absent from mathlib; we provide a complete construction with all seven \texttt{StarAlgEquiv} field proofs, using only associativity and unitarity.
  \item \texttt{reindexStarAlgEquiv} --- upgrades \texttt{Matrix.reindexAlgEquiv} to preserve $\star$ (via \texttt{Matrix.conjTranspose\_reindex}, which is definitionally \texttt{rfl} when the same reindex equivalence is applied to both row and column axes).
  \item \texttt{coe\_iSup\_of\_directed\_starSubalgebra} --- mathlib provides the corresponding lemma for \texttt{NonUnitalStarSubalgebra} but not for the unital version; we fill the gap in $\sim$25 lines.
  \item A $C^{\ast}$-algebraic direct-limit construction for the non-commutative ternary tower ($\texttt{TimelessFieldRing}$ and its equipment) --- mathlib's \texttt{Ring.DirectLimit} requires \texttt{CommRing} and does not adapt; the generic \texttt{DirectLimit} lacks metric structure.
  \item A star lift for \texttt{UniformSpace.Completion} --- mathlib has \texttt{UniformSpace.Completion.extensionHom} at the \texttt{RingHom} level but not at the $\ast$-algebra level. Our construction $\texttt{TimelessFieldRingToCompletionStarAlgHom}$ specializes the pattern and could be generalized to a mathlib helper $\texttt{IsDenseInducing.extendStarAlgHom}$.
  \item The $\texttt{IsTracialLinearFunctional}$ predicate itself, with an associated uniqueness theorem for the $3^{\infty}$ substrate --- mathlib has zero coverage of tracial states.
\end{enumerate}

Each of the first two are less than 50 lines and could be submitted as standalone mathlib PRs. Items 3--6 would require broader discussion with the mathlib community regarding scope and naming conventions.

\section{Kernel-audit verdict}
\label{sec:audit}

Running \texttt{lake build PF.SubstrateRigidity} against the current \texttt{mathlib4} pin (\texttt{v4.24.0-rc1}) yields the following audit block, output by the six \texttt{\#print axioms} directives at the tail of the file:

\begin{lstlisting}[language={}]
T_infinity_rigidity                    : [propext, Classical.choice, Quot.sound]
substrate3Inf_iso                      : [propext, Classical.choice, Quot.sound]
substrate3Inf_TimelessFieldCompletion  : [propext, Classical.choice, Quot.sound]
Substrate3Inf.connect_iso              : [propext, Classical.choice, Quot.sound]
Substrate3Inf.connect_unital           : [propext, Classical.choice, Quot.sound]
IsTracialLinearFunctional              : [propext, Classical.choice, Quot.sound]

Build completed successfully (3265 jobs).
\end{lstlisting}

\noindent
No \texttt{sorryAx}, no \texttt{ofReduceBool}, no project axioms. The three axioms shown are the standard Lean kernel axioms (propositional extensionality, classical choice, and quotient soundness), assumed by \texttt{mathlib4} throughout.

\section{Place in Principia Fractalis}
\label{sec:pf-context}

This result should not be read in isolation from the research program in which the substrate object $T_{\infty}$ was constructed. \emph{Principia Fractalis}~\cite{cohen_pf_2026} is Pablo Cohen's long-running (2024--present) research program, presenting a candidate substrate-level framework based on a \emph{Fractal Resonance Function}
\[
R_f(\alpha, s) \;=\; \sum_{n \geq 1} \frac{e^{i \pi \alpha \, D_3(n)}}{n^s}
\]
(where $D_3$ is the base-3 digit sum), with the Timeless Field $T_{\infty}$ --- the object formalized here as \texttt{TimelessFieldCompletion} --- serving as its algebraic substrate. The program claims that a nine-parameter $\alpha$-skeleton, subject to eight structural laws, projects into a family of physical and mathematical sectors, and it advances candidate answers to several of the Clay Millennium Problems~\cite{clay_problems}.

The Principia Fractalis codebase, hosted at \texttt{github.com/FractalDevTeam/Principia-Fractalis}, comprises approximately $4.3 \times 10^5$ lines of Lean~4 with zero \texttt{sorry} declarations in the main build (excluding the \texttt{PF.SubstrateRigidity} module presented here, which was not yet wired into the top-level import when this paper was written). It includes over $10^4$ machine-verified theorems, a book manuscript of 36 chapters and 13 appendices, and 15 accompanying research papers.

\subsection*{Where the joint statement lands in the program}

The paired result (Corollary~\ref{cor:joint}) discharges the ``central theorem'' specified in \S 1 of the program's own steering document \texttt{UNIFIED\_THEORY\_PROOF\_PROGRAM.md} (dated 2026-09-09):

\begin{quote}\itshape
The Timeless Field substrate $T_{\infty}$ is the UHF algebra of supernatural type $3^{\infty}$ [\ldots]. Any $C^{\ast}$-algebra satisfying the substrate axioms (ternary directed system, norm-preserving connecting maps, $C^{\ast}$-identity) is $\ast$-isomorphic to it.
\end{quote}

\noindent
Theorem~\ref{thm:main} is that discharge, formalized. In the language of Stone--von~Neumann or GNS it is a \emph{rigidity theorem}: the substrate is not a modelling choice --- it is the unique object of its class, up to $\ast$-isomorphism, characterized by four Lean fields.

The pairing with Theorem~\ref{thm:obstruction} sharpens the reading and, in our view, matches what the program has actually established about its substrate axis. What follows is the reading we recommend for downstream discussion:

\subsection*{The two-axis picture}

Principia Fractalis has \emph{two} independent axes at its foundation:

\begin{description}[leftmargin=2em]
  \item[Axis 1 --- the substrate.] The $C^{\ast}$-algebraic object $T_{\infty}$. Kernel-verified as of the present work to be uniquely determined by its ternary UHF structure and its unique tracial linear functional. This paper's contribution.
  \item[Axis 2 --- the $\alpha$-skeleton.] The nine real-valued parameters $\alpha_{\mathrm{Poincar\acute{e}}} = 1$, $\alpha_{\mathrm{Hodge}} = \varphi$, $\alpha_{\mathrm{RH}} = 3/2$, $\alpha_{\mathrm{YM}} = 2$, $\alpha_{\mathrm{BSD}} = 3\pi/4$, $\alpha_{\mathrm{NS}} = 3\pi/2$, $\alpha_{\mathrm{P}} = \sqrt{2}$, $\alpha_{\mathrm{NP}} = \varphi + 1/4$, $\alpha_{\mathrm{QG}} = \sqrt{2\pi}$. Kernel-verified in the parent project (Cohen, r128) to be uniquely determined \emph{given} eight structural laws and the Perelman anchor $\alpha_{\mathrm{Poincar\acute{e}}} = 1$. The eight laws are, per the rigidity audit~\cite{cohen_alpha_rigidity_2026}, \emph{independent} assumptions.
\end{description}

Theorem~\ref{thm:obstruction} establishes that the two axes are not linked through the substrate's $K$-theoretic channel: seven of nine $\alpha$-values are irrational (and one further ratio, $\alpha_{\mathrm{NS}}/\alpha_{\mathrm{RH}}$, equals $\pi$), placing them outside $\mathbb{Z}[1/3]$; the substrate's classifying invariant has image exactly $\mathbb{Z}[1/3]$. Any bridge between the two axes must therefore be built at a higher layer than the substrate --- through, for instance, sector-specific interface theorems (per the program's own directive, \S 5).

\subsection*{What Theorem~\ref{thm:main} \emph{does not} entail}

Directly and without hedging:

\begin{itemize}[leftmargin=*]
  \item \textbf{No $\alpha$-value is derived by the substrate.} Theorem~\ref{thm:obstruction} makes this obstruction explicit through the $K$-theoretic channel; the rigidity audit~\cite{cohen_alpha_rigidity_2026} makes it explicit through the direct-membership channel; between them, the ``$\alpha$-values as substrate consequences'' reading is closed. The nine $\alpha$-values remain, in the audit's own recommended phrasing, ``a presentation of the nine values, not a derivation of them.''
  \item \textbf{No Clay Millennium Problem is resolved.} The rigidity theorem does not touch the Riemann Hypothesis, Yang--Mills mass gap, Navier--Stokes regularity, BSD, Hodge, or P versus NP. The parent program's own machine-verified negative results (chapter-level refutations recorded in the codex tree; see~\cite{cohen_pf_2026}) are unaffected either way.
  \item \textbf{No sector correspondence is established.} The program's Section~5 obligation (per-sector interface theorems bridging the substrate to standard mathematical objects) is untouched by this result. That work is ongoing; the substrate layer is now stable enough to serve as an anchor for it, but the interfaces themselves must be built.
\end{itemize}

\subsection*{What the joint statement is worth}

We claim, narrowly and, we believe, defensibly:
\begin{enumerate}[leftmargin=2em]
  \item Principia Fractalis's substrate is a well-defined, unique-up-to-isomorphism $C^{\ast}$-algebraic object, kernel-certified.
  \item The substrate's classifying invariant provably underdetermines the framework's $\alpha$-skeleton, and any derivation of the $\alpha$-values must be established at a higher layer.
  \item Whatever downstream mathematical or physical content the program eventually establishes will rest on a substrate that is not a modelling degree of freedom, and on an $\alpha$-skeleton whose grounding is a subject for continued research.
\end{enumerate}

We regard this as the honest scientific summary of what has been kernel-verified about the substrate axis of Principia Fractalis through 2026-09-10. It is neither a headline claim about the framework's Clay-adjacent aspirations, nor a dismissal of them --- it is the substrate-side theorem plus the substrate-side obstruction, offered as a joint statement whose downstream implications remain for the sector-interface work to determine.

\section{Related work}
\label{sec:related}

\paragraph{Formalization of C*-algebras.} Buzzard, Nash, Riccardo, and others have contributed substantially to \texttt{mathlib4}'s $C^{\ast}$-algebra infrastructure, culminating in the current \texttt{CStarAlgebra} class hierarchy~\cite{mathlib_cstar_2026}. Our work builds on this infrastructure and, where it falls short, fills the gaps in a manner consistent with mathlib conventions. To our knowledge, no prior Lean formalization has treated inductive limits of matrix algebras in a $C^{\ast}$-algebraic (as opposed to purely algebraic) manner.

\paragraph{Formalization of classification theorems.} Alper, Franke, and others have formalized the classification of finite simple groups in various fragmentary states; Kepler's conjecture~\cite{hales_kepler_formal_2017} and the four-color theorem~\cite{gonthier_4color_2008} are the standard classification-adjacent success stories. Most recently, Anthropic's formalization of Fermat's Last Theorem~\cite{anthropic_flt_2026} demonstrated the practical viability of large-scale language-model-assisted formalization, from which the present work adopts the statement-card discipline.

\paragraph{Bratteli diagrams and Elliott invariants.} A full formalization of Elliott's classification for AF~$C^{\ast}$-algebras~\cite{elliott1976} would be a much more ambitious undertaking than the present specialization; we regard the present work as an initial move in that direction, with the UHF $3^{\infty}$ specialization serving as a proof-of-concept.

\section{Future work}
\label{sec:future}

The formalization is a specialization; the natural generalizations are:
\begin{enumerate}[leftmargin=*]
  \item UHF algebras of arbitrary supernatural number $\prod_p p^{n_p}$. Requires formalizing supernatural numbers themselves (a natural type constructor), Bratteli diagrams, and re-parameterizing the tower.
  \item AF ($C^{\ast}$)-algebras --- Elliott's full 1976 classification via the ordered $K_0$ group. Substantially more machinery; a plausible medium-term target.
  \item Extending \texttt{IsTracialLinearFunctional} to include positivity and star-hermitianness --- currently omitted (see Remark~\ref{rem:tracial-caveat}) but readily added. Would tighten the semantic content of the trace hypothesis to match the standard definition.
  \item Building the sector-interface theorems in the parent Principia Fractalis program on top of the now-stable substrate layer (per the program's own directive, \S 5).
\end{enumerate}

\section{Reproducibility and availability}
\label{sec:availability}

The formalization is a single file, \texttt{PF/SubstrateRigidity.lean}, of approximately 2{,}400 lines, on the \texttt{r331b-provenance} branch of the public repository \texttt{FractalDevTeam/Principia-Fractalis} at commit \texttt{18f55a14} (2026-09-10). To reproduce the audit verdict:

\begin{lstlisting}[language={}]
git clone git@github.com:FractalDevTeam/Principia-Fractalis.git
cd Principia-Fractalis
git checkout r331b-provenance
cd PF_Lean4_Code
lake build PF.SubstrateRigidity
\end{lstlisting}

\noindent
The pinned toolchain is \texttt{leanprover/lean4:v4.24.0-rc1} with \texttt{mathlib4} at its \texttt{lakefile.lean}-pinned version. Build completes in approximately 8 minutes on stock hardware (Ryzen 9~7950X, 32~GB~RAM).

\section*{Acknowledgments}

The first author acknowledges the ORCID identifier \texttt{0009-0002-0734-5565} and thanks the mathlib community for the deep $C^{\ast}$-algebra infrastructure that made this work possible. Both authors thank Anthropic for the compute infrastructure and the Claude Code / Claude Agents platform on which the coordinated proof-farming discipline was implemented, adopting patterns from Anthropic's own Fermat's Last Theorem formalization~\cite{anthropic_flt_2026}. All mathematical decisions --- theorem specifications, proof strategy, semantic scope --- were the first author's; the Lean 4 proof engineering was farmed via prover-agent sessions with human-authored statement cards and human-audited verdicts.

\begin{thebibliography}{99}

\bibitem{glimm1960}
J.~G. Glimm.
\newblock On a certain class of operator algebras.
\newblock \emph{Trans. Amer. Math. Soc.} 95:2, 318--340, 1960.

\bibitem{elliott1976}
G.~A. Elliott.
\newblock On the classification of inductive limits of sequences of semisimple finite-dimensional algebras.
\newblock \emph{J. Algebra} 38, 29--44, 1976.

\bibitem{skolem1927}
T.~Skolem.
\newblock Zur Theorie der assoziativen Zahlensysteme.
\newblock \emph{Skrifter Norsk Vid. Akad. Oslo} I, 12, 1927.

\bibitem{noether1929}
E.~Noether.
\newblock Hyperkomplexe Gr\"o\ss{}en und Darstellungstheorie.
\newblock \emph{Math. Z.} 30, 641--692, 1929.

\bibitem{clay_problems}
Clay Mathematics Institute.
\newblock Millennium Prize Problems.
\newblock \url{https://www.claymath.org/millennium-problems/}, accessed 2026-09-10.

\bibitem{cohen_pf_2026}
P.~Cohen.
\newblock \emph{Principia Fractalis} (in preparation), 2024--2026.
\newblock Manuscript and formalization repository: \url{https://github.com/FractalDevTeam/Principia-Fractalis}.

\bibitem{cohen_alpha_rigidity_2026}
P.~Cohen.
\newblock $\alpha$-skeleton rigidity audit --- adversarial report.
\newblock Codex record dated 2026-09-07, Principia Fractalis repository, file \texttt{codex/ALPHA\_RIGIDITY\_AUDIT\_REPORT\_2026-09-07.md}.
\newblock Kernel artifacts: \texttt{PF/AlphaL5PiScalingObstruction\_r332.lean}, \texttt{PF/AlphaStructuralLawAudit\_r334.lean}, \texttt{PF/AlphaStructuralLawAuditI7\_r335.lean}.

\bibitem{mathlib_cstar_2026}
The mathlib4 Community.
\newblock \emph{Mathlib4: The Lean 4 Mathematical Library}, C*-algebra namespace.
\newblock Continuously updated; snapshot at commit HEAD, 2026-09-09.
\newblock \url{https://leanprover-community.github.io/mathlib4_docs/Mathlib/Analysis/CStarAlgebra/}.

\bibitem{anthropic_flt_2026}
Anthropic.
\newblock A formalization of Fermat's Last Theorem in Lean 4.
\newblock Research report, September 2026.
\newblock \url{https://www.anthropic.com/research/}.

\bibitem{hales_kepler_formal_2017}
T.~C. Hales et al.
\newblock A formal proof of the Kepler conjecture.
\newblock \emph{Forum of Mathematics, Pi} 5, e2, 2017.

\bibitem{gonthier_4color_2008}
G.~Gonthier.
\newblock Formal proof --- the four-color theorem.
\newblock \emph{Notices of the AMS} 55:11, 1382--1393, 2008.

\end{thebibliography}

\appendix

\section{Full Lean signatures of the load-bearing declarations}
\label{app:signatures}

For reference, the exact Lean 4 signatures of the six kernel-audited declarations:

\begin{lstlisting}
-- 1. The tracial linear functional predicate.
structure IsTracialLinearFunctional
    (A : Type*) [CStarAlgebra A] (φ : A → ℂ) : Prop where
  continuous : Continuous φ
  add        : ∀ x y : A, φ (x + y) = φ x + φ y
  smul       : ∀ (c : ℂ) (x : A), φ (c • x) = c * φ x
  tracial    : ∀ x y : A, φ (x * y) = φ (y * x)
  unital     : φ 1 = 1

-- 2. The characterization structure.
structure Substrate3Inf (A : Type*) [CStarAlgebra A] where
  tower         : ℕ → StarSubalgebra ℂ A
  tower_matrix  : ∀ k, Nonempty
                    (tower k ≃⋆ₐ[ℂ] Matrix (Fin (3^k)) (Fin (3^k)) ℂ)
  tower_mono    : ∀ k, tower k ≤ tower (k+1)
  tower_dense   : Dense (((⨆ k, tower k : StarSubalgebra ℂ A) : Set A))
  trace_unique  : ∃! τ : A → ℂ, IsTracialLinearFunctional A τ

-- 3. Auxiliary lemmas on the structure.
lemma Substrate3Inf.connect_unital
    {A : Type*} [CStarAlgebra A] (h : Substrate3Inf A) (k : ℕ) :
    (StarSubalgebra.inclusion (h.tower_mono k)) 1 = 1

lemma Substrate3Inf.connect_iso
    {A : Type*} [CStarAlgebra A] (h : Substrate3Inf A) (k : ℕ) :
    Isometry (StarSubalgebra.inclusion (h.tower_mono k))

-- 4. The general classification theorem.
theorem substrate3Inf_iso
    {A B : Type*} [CStarAlgebra A] [CStarAlgebra B]
    (hA : Substrate3Inf A) (hB : Substrate3Inf B) :
    Nonempty (A ≃⋆ₐ[ℂ] B)

-- 5. The witness on the concrete substrate.
noncomputable def substrate3Inf_TimelessFieldCompletion :
    Substrate3Inf TimelessFieldCompletion

-- 6. The main theorem.
theorem T_infinity_rigidity
    (A : Type*) [CStarAlgebra A]
    (h : Substrate3Inf A) :
    Nonempty (A ≃⋆ₐ[ℂ] TimelessFieldCompletion) :=
  substrate3Inf_iso h substrate3Inf_TimelessFieldCompletion
\end{lstlisting}

\end{document}
